\documentclass[11pt,a4paper]{article}
\usepackage{shared/luxcover}
\usepackage[utf8]{inputenc}
\usepackage[T1]{fontenc}
\usepackage{amsmath,amssymb,amsfonts,amsthm}
\usepackage{graphicx}
\usepackage{hyperref}
\usepackage{algorithm}
\usepackage{algpseudocode}
\usepackage{booktabs}
\usepackage{enumitem}
\usepackage{tikz}
\usetikzlibrary{arrows.meta,positioning,shapes.geometric}

\newtheorem{theorem}{Theorem}[section]
\newtheorem{lemma}[theorem]{Lemma}
\newtheorem{proposition}[theorem]{Proposition}
\newtheorem{corollary}[theorem]{Corollary}
\newtheorem{definition}{Definition}[section]
\newtheorem{remark}{Remark}[section]

\hypersetup{
    colorlinks=true,
    linkcolor=black,
    citecolor=blue,
    urlcolor=blue
}

\title{\textbf{Lux as Trust Kernel:\\Identity, Key Policy, Threshold Governance,\\and Receipts}}

\author{
Zach Kelling\thanks{Corresponding author: zach@lux.network}\\
\textit{Lux Industries}\\
\texttt{research@lux.network}
}

\date{2026}

\begin{document}

\luxcoverpage

\begin{abstract}
We formalize Lux Network as a \emph{trust kernel}: a minimal on-chain substrate that
provides identity binding, key lifecycle policy, threshold governance, and tamper-evident
receipts---without storing application data on-chain. The trust kernel model inverts the
common blockchain-as-database pattern. Applications retain full data sovereignty in local
encrypted vaults (capsules) while delegating exactly four concerns to the chain:
(i)~globally-unique identity via decentralized identifiers (\texttt{did:lux:*}),
(ii)~key registration and rotation policy on the K-Chain,
(iii)~quorum-controlled reveal and treasury approval on the T-Chain with TFHE policy
evaluation and CKKS confidential compute,
and (iv)~multi-party threshold signatures on the M-Chain via FROST and CGGMP21.
Every state transition produces a cryptographic receipt anchored on-chain, enabling
provider-independent auditability. We define the trust kernel formally, prove that it
preserves capsule sovereignty under Byzantine fault tolerance assumptions, and show
that the architecture enables cross-capsule coordination (key exchange, capability
delegation, threshold approval) without revealing plaintext to any single party.
All critical protocols are verified in Lean~4 against mechanized proofs for
consensus safety, FROST unforgeability, CGGMP21 threshold composition, and
trust authority delegation.

\textbf{Keywords}: trust kernel, decentralized identity, threshold governance,
key policy, receipt anchoring, post-quantum cryptography, TFHE, CKKS, FROST, CGGMP21
\end{abstract}

\tableofcontents
\newpage

%% ===================================================================
\section{Introduction}\label{sec:intro}
%% ===================================================================

Blockchains are routinely misapplied as general-purpose databases. This conflation
creates systems that are simultaneously too expensive for data storage and too
opaque for trust verification. The root cause is a failure to distinguish between
two orthogonal concerns: \emph{data management} (storage, indexing, query) and
\emph{trust management} (identity, authorization, coordination, audit).

\begin{definition}[Trust Kernel]\label{def:trust-kernel}
A trust kernel $\mathcal{K}$ is a replicated state machine that provides exactly
four services to application capsules:
\begin{enumerate}[label=(\alph*)]
    \item \textbf{Identity}: globally-unique, self-sovereign identifiers bound to
          cryptographic key material;
    \item \textbf{Key Policy}: lifecycle management (registration, rotation,
          revocation) with enforceable policy constraints;
    \item \textbf{Threshold Governance}: quorum-controlled state transitions
          requiring $t$-of-$n$ authorization;
    \item \textbf{Receipts}: tamper-evident, append-only records of every
          trust-relevant state transition.
\end{enumerate}
\end{definition}

The trust kernel stores no application data. It stores \emph{commitments to decisions
about data}: who created it, who may access it, under what conditions, and when those
conditions were evaluated. This separation is not merely architectural preference---it
is a security boundary. Data that never touches the chain cannot leak from the chain.

\subsection{What Local Vaults Cannot Provide}

An encrypted local vault (a capsule, as defined in \cite{lux-capsule-architecture})
provides strong confidentiality and integrity for data at rest. However, a vault alone
cannot provide:

\begin{enumerate}[label=(\roman*)]
    \item \textbf{Global identity}: a vault can generate keypairs, but cannot prove
          uniqueness or non-repudiation to a third party without a shared root of trust;
    \item \textbf{Cross-party coordination}: two vaults cannot agree on a shared secret,
          delegate capabilities, or approve a joint action without a coordination layer
          that both trust;
    \item \textbf{Settlement finality}: a vault can record local state, but cannot
          prove to an external verifier that a state transition occurred at a specific
          time and was not subsequently reversed.
\end{enumerate}

The trust kernel fills exactly these gaps. It is the minimum viable chain: the smallest
on-chain footprint that makes sovereign applications composable.

\subsection{Contribution}

This paper makes the following contributions:

\begin{enumerate}
    \item We formalize the trust kernel abstraction and prove it sufficient for
          sovereign application composition (Section~\ref{sec:intro}).
    \item We specify the K-Chain key registry with ML-KEM wrapping, rotation
          policy, and revocation receipts (Section~\ref{sec:kchain}).
    \item We define T-Chain threshold governance with TFHE policy evaluation and
          CKKS confidential compute (Section~\ref{sec:tchain}).
    \item We describe M-Chain MPC signing via FROST and CGGMP21
          (Section~\ref{sec:mchain}).
    \item We specify the I-Chain decentralized identity layer with ZK selective
          disclosure (Section~\ref{sec:ichain}).
    \item We define the receipt anchoring protocol and prove its
          tamper-evidence (Section~\ref{sec:receipts}).
    \item We show how the trust kernel enables provider switching without data
          loss (Section~\ref{sec:provider}).
    \item We formalize cross-capsule coordination through the trust kernel
          (Section~\ref{sec:cross-capsule}).
    \item We present the security model with formal trust assumptions and PQ
          migration path (Section~\ref{sec:security}).
    \item We reference mechanized Lean~4 proofs for all critical protocols
          (Section~\ref{sec:verification}).
\end{enumerate}

%% ===================================================================
\section{K-Chain: Key Registry and Policy}\label{sec:kchain}
%% ===================================================================

The K-Chain is a purpose-built chain within the Lux primary network that manages
the lifecycle of cryptographic keys. It does not store encrypted data; it stores
\emph{policy about keys}.

\subsection{Key Registry}

\begin{definition}[Key Record]\label{def:key-record}
A key record $\kappa$ is a tuple:
\[
    \kappa = (\mathsf{id}, \mathsf{pk}, \mathsf{alg}, \mathsf{caps}, \mathsf{policy},
              \mathsf{created}, \mathsf{expires}, \mathsf{status})
\]
where $\mathsf{id} \in \{0,1\}^{256}$ is the key identifier, $\mathsf{pk}$ is the
public key, $\mathsf{alg} \in \{\text{ML-KEM-768}, \text{ML-KEM-1024}, \text{Ed25519},
\text{ML-DSA-65}, \text{secp256k1}\}$ is the algorithm identifier, $\mathsf{caps}$
is a capability bitfield, $\mathsf{policy}$ is a policy predicate,
$\mathsf{created}$ and $\mathsf{expires}$ are block heights, and
$\mathsf{status} \in \{\textsc{active}, \textsc{rotated}, \textsc{revoked}\}$.
\end{definition}

\subsection{ML-KEM Wrapping}

Keys stored in capsule vaults are wrapped using ML-KEM (FIPS~203) as specified in
\cite{lux-pq-crypto-suite}. The K-Chain stores only the public encapsulation key
$\mathsf{ek}$ and the policy governing who may request encapsulation. The wrapped
key material never appears on-chain.

\begin{definition}[Wrapping Policy]\label{def:wrap-policy}
A wrapping policy $\pi_w$ for key record $\kappa$ is a predicate over requestor
identity and context:
\[
    \pi_w : \mathsf{DID} \times \mathsf{Context} \to \{\top, \bot\}
\]
where $\mathsf{Context}$ includes the requestor's capability set, the current block
height, and an optional threshold requirement.
\end{definition}

When a capsule requires a wrapped key, it submits a \textsc{WrapRequest} transaction
to the K-Chain. The chain evaluates $\pi_w$ against the requestor's DID and, if
satisfied, emits a \textsc{WrapGrant} receipt containing the ML-KEM ciphertext
$c = \mathsf{Encaps}(\mathsf{ek}; r)$ and the shared secret $K$ is delivered
to the requestor off-chain via a secure channel.

\subsection{Rotation Policy}

\begin{definition}[Rotation Rule]\label{def:rotation}
A rotation rule $\rho$ for key $\kappa$ specifies:
\[
    \rho = (\mathsf{max\_age}, \mathsf{successor\_alg}, \mathsf{approval\_threshold},
            \mathsf{grace\_period})
\]
The K-Chain enforces that no key remains active beyond $\mathsf{max\_age}$ blocks
past $\kappa.\mathsf{created}$. Upon rotation, the old key enters
$\textsc{rotated}$ status and a new key record is created with algorithm
$\mathsf{successor\_alg}$.
\end{definition}

Rotation may be triggered by three events: (1)~age expiry, (2)~explicit owner request,
or (3)~threshold governance vote on the T-Chain. In cases (1) and (3), the rotation
is mandatory and produces a \textsc{RotationReceipt} regardless of owner action.

\subsection{Capability Grants and Revocation}

Capabilities are delegated via \textsc{GrantCap} transactions that append to the
key record's capability set. Each grant is scoped to a specific DID, a capability
type (read, sign, admin), and an optional time bound.

Revocation is immediate and final. A \textsc{RevokeCap} transaction sets the
capability's status to revoked and emits a \textsc{RevocationReceipt}. The receipt
includes the revoked capability hash, the revoking authority's DID, and the block
height, providing a permanent audit trail.

\begin{proposition}[Revocation Finality]\label{prop:revocation}
Once a \textsc{RevocationReceipt} is finalized on the K-Chain, no subsequent
transaction can restore the revoked capability for the same key--DID pair.
\end{proposition}

%% ===================================================================
\section{T-Chain: Threshold Governance}\label{sec:tchain}
%% ===================================================================

The T-Chain handles quorum-controlled operations: actions that require $t$-of-$n$
authorization before execution. Unlike simple multi-signature schemes, the T-Chain
supports \emph{policy evaluation over encrypted inputs} via TFHE and \emph{confidential
aggregation} via CKKS.

\subsection{Quorum-Controlled Reveal}

\begin{definition}[Quorum Policy]\label{def:quorum}
A quorum policy $Q$ is a tuple:
\[
    Q = (n, t, \mathsf{voters}, \mathsf{predicate}, \mathsf{timeout})
\]
where $n$ is the total number of authorized voters, $t$ is the threshold,
$\mathsf{voters} \subseteq \mathsf{DID}^n$ is the voter set,
$\mathsf{predicate}$ is a boolean circuit evaluated on voter inputs,
and $\mathsf{timeout}$ is a block height deadline.
\end{definition}

Applications of quorum-controlled reveal include:
\begin{itemize}
    \item \textbf{Treasury approval}: $t$-of-$n$ board members must approve a
          disbursement before the transaction is signed on M-Chain.
    \item \textbf{Admin recovery}: a compromised admin key can be rotated if $t$
          recovery contacts authorize the rotation on K-Chain.
    \item \textbf{Conditional access}: a document is revealed to a requestor only
          if $t$ data stewards approve the access request.
\end{itemize}

\subsection{TFHE Policy Evaluation}

For policies where voter inputs must remain private (e.g., salary-based approval
thresholds, credit scoring), the T-Chain evaluates the quorum predicate as a
boolean circuit over TFHE-encrypted bits. Each voter submits an encrypted vote
$\hat{v}_i = \mathsf{TFHE.Enc}(v_i)$. The chain evaluates the predicate
circuit homomorphically:
\[
    \hat{r} = \mathsf{Circuit}(\hat{v}_1, \ldots, \hat{v}_n)
\]
and the result is threshold-decrypted by the T-Chain validators, yielding
$r \in \{0, 1\}$ without revealing individual votes. The TFHE parameters
follow \cite{lux-tfhe}: $N = 1024$, $q \approx 2^{27}$, 128-bit security.

\subsection{CKKS Confidential Compute}

For policies requiring arithmetic over real-valued inputs (e.g., weighted voting,
aggregate risk scores), the T-Chain supports CKKS approximate homomorphic encryption.
Each party submits $\hat{x}_i = \mathsf{CKKS.Enc}(x_i)$ where $x_i \in \mathbb{R}$.
The chain computes:
\[
    \hat{y} = \sum_{i=1}^{n} w_i \cdot \hat{x}_i
\]
and compares the result against a threshold $\tau$. The aggregated value $y$ is
revealed; individual inputs $x_i$ are not.

\begin{proposition}[Vote Privacy]\label{prop:vote-privacy}
Under the LWE assumption with parameters $(N, q, \sigma)$ as specified in
\cite{lux-pq-crypto-suite}, no coalition of fewer than $t$ T-Chain validators
can recover any individual vote $v_i$ from the transcript of a quorum evaluation.
\end{proposition}

\subsection{Treasury Approval and Admin Recovery}

Treasury operations bind T-Chain governance to M-Chain signing. A disbursement
follows a strict sequence:

\begin{enumerate}
    \item A proposer submits a \textsc{TreasuryProposal} to the T-Chain with
          the target address, amount, and justification hash.
    \item Authorized voters submit encrypted approvals.
    \item The T-Chain evaluates the quorum predicate.
    \item If approved, the T-Chain emits a \textsc{GovernanceReceipt} that
          authorizes the M-Chain to produce a threshold signature on the
          disbursement transaction.
    \item The M-Chain produces the signature only upon verifying the receipt.
\end{enumerate}

Admin recovery follows the same pattern with the rotation request replacing the
disbursement transaction. The recovery quorum is stored on the K-Chain as part
of the key record's policy (Definition~\ref{def:rotation}).

%% ===================================================================
\section{M-Chain: MPC Signing}\label{sec:mchain}
%% ===================================================================

The M-Chain provides threshold signature generation as a chain-native service. It
implements two protocols: FROST for Schnorr-based curves (Ed25519, BIP-340) and
CGGMP21 for ECDSA (secp256k1). The full specification is given in
\cite{lux-mchain-mpc}; here we describe the trust kernel interface.

\subsection{FROST Threshold Signatures}

\begin{definition}[FROST Signing Session]\label{def:frost}
A FROST signing session $\mathcal{S}$ for message $m$ under group public key
$\mathsf{gpk}$ with threshold $t$ proceeds in two rounds:
\begin{enumerate}
    \item \textbf{Commitment}: each signer $i$ generates nonce pair
          $(d_i, e_i) \gets \mathbb{Z}_q^2$ and broadcasts commitments
          $(D_i = g^{d_i}, E_i = g^{e_i})$.
    \item \textbf{Response}: each signer $i$ computes partial signature
          $z_i = d_i + \rho_i \cdot e_i + \lambda_i \cdot s_i \cdot c$
          where $c = H(\mathsf{gpk}, R, m)$, $\rho_i$ is the binding factor,
          and $\lambda_i$ is the Lagrange coefficient.
\end{enumerate}
The aggregator computes $\sigma = (R, z)$ where $z = \sum_{i \in S} z_i$ and $S$
is the signer set with $|S| \geq t$.
\end{definition}

\subsection{CGGMP21 Threshold ECDSA}

For secp256k1 operations (Ethereum, Bitcoin legacy), the M-Chain uses the CGGMP21
protocol \cite{lux-mchain-mpc} which provides:
\begin{itemize}
    \item Pre-signing phase with Paillier commitments for multiplicative-to-additive
          share conversion.
    \item Identifiable abort: if a signer submits an invalid partial signature, the
          protocol identifies and slashes the faulty party.
    \item Non-interactive online phase after a single pre-signing round.
\end{itemize}

\subsection{Authorization Gate}

The M-Chain will not produce a signature without a valid authorization. Authorization
sources are:

\begin{enumerate}
    \item A \textsc{GovernanceReceipt} from the T-Chain (for treasury operations).
    \item A \textsc{WrapGrant} from the K-Chain (for key operations).
    \item A valid DID authentication proof from the I-Chain (for user-initiated
          transactions).
\end{enumerate}

\begin{theorem}[Signing Authorization]\label{thm:signing-auth}
Let $\sigma$ be a threshold signature produced by the M-Chain. Then there exists
a finalized receipt $\mathcal{R}$ on at least one of \{K-Chain, T-Chain, I-Chain\}
that authorized the signing session. No valid signature can be produced without
such a receipt.
\end{theorem}

\begin{proof}
The M-Chain VM accepts \textsc{SignRequest} transactions only if accompanied by
a receipt reference $(chain, block, txIndex)$. The VM verifies the receipt via
Warp messaging \cite{lux-warp-messaging} before initiating the signing protocol.
A \textsc{SignRequest} without a valid receipt reference is rejected at the
transaction validation layer, before entering the mempool. Since the signing
protocol requires the \textsc{SignRequest} to be finalized on the M-Chain,
and finalization requires valid receipt verification, the theorem holds.
\end{proof}

%% ===================================================================
\section{I-Chain: Decentralized Identity}\label{sec:ichain}
%% ===================================================================

The I-Chain implements the \texttt{did:lux:*} method as specified in
\cite{lux-id-did-specification}. Within the trust kernel, the I-Chain serves
as the root of all identity assertions.

\subsection{DID Document Structure}

\begin{definition}[Lux DID]\label{def:did}
A Lux DID has the form \texttt{did:lux:<method-specific-id>} where the
method-specific identifier is derived from the subject's initial public key:
\[
    \mathsf{msid} = \mathsf{Base58}(\mathsf{RIPEMD160}(\mathsf{SHA256}(\mathsf{pk}_0)))
\]
The DID document contains verification methods (public keys), authentication
relationships, service endpoints, and a controller field.
\end{definition}

Each DID document is registered on the I-Chain as an immutable record. Updates
produce new versions linked to the previous version by hash, forming a verifiable
history chain.

\subsection{Verifiable Credentials}

The I-Chain supports W3C Verifiable Credentials with Lux-specific extensions:
\begin{itemize}
    \item \textbf{Credential issuance}: an issuer signs a credential binding a
          claim to a subject DID and registers the credential hash on the I-Chain.
    \item \textbf{Credential verification}: a verifier checks the issuer's
          signature and confirms the credential hash exists on-chain with status
          $\textsc{active}$.
    \item \textbf{Credential revocation}: the issuer submits a
          \textsc{RevokeCredential} transaction; the credential's on-chain status
          changes to $\textsc{revoked}$.
\end{itemize}

\subsection{ZK Selective Disclosure}

Credentials may be presented with zero-knowledge selective disclosure. The holder
generates a ZK proof that demonstrates possession of a credential satisfying a
predicate (e.g., ``age $\geq$ 18'') without revealing the full credential.

\begin{definition}[Selective Disclosure Proof]\label{def:zksd}
Given a credential $C$ with attributes $\{a_1, \ldots, a_k\}$ and a predicate
$\phi$ over a subset of attributes, the holder produces a proof $\pi$ such that:
\[
    \mathsf{Verify}(\pi, \phi, \mathsf{issuer\_pk}, \mathsf{holder\_did}) = \top
    \iff \phi(a_{i_1}, \ldots, a_{i_m}) = \top
\]
where $\{i_1, \ldots, i_m\} \subseteq \{1, \ldots, k\}$ and no information
about attributes outside the predicate is revealed.
\end{definition}

\subsection{Credential Revocation via Accumulator}

The I-Chain maintains a cryptographic accumulator for credential revocation.
Adding a credential hash to the revocation accumulator constitutes revocation.
Verification requires a non-membership witness: the holder proves their credential
is \emph{not} in the revocation set without downloading the full revocation list.

%% ===================================================================
\section{Receipt Anchoring}\label{sec:receipts}
%% ===================================================================

Every trust-relevant state transition in the Lux trust kernel produces a
\emph{receipt}: a compact, self-contained proof that the transition occurred.

\begin{definition}[Trust Receipt]\label{def:receipt}
A trust receipt $\mathcal{R}$ is a tuple:
\[
    \mathcal{R} = (\mathsf{chain}, \mathsf{block}, \mathsf{txHash},
                   \mathsf{type}, \mathsf{subject}, \mathsf{digest},
                   \mathsf{sig})
\]
where $\mathsf{chain} \in \{\text{K}, \text{T}, \text{M}, \text{I}\}$ identifies
the originating chain, $\mathsf{block}$ and $\mathsf{txHash}$ identify the
transaction, $\mathsf{type}$ is the receipt type (Table~\ref{tab:receipts}),
$\mathsf{subject}$ is the affected DID, $\mathsf{digest}$ is the SHA-256 hash
of the receipt payload, and $\mathsf{sig}$ is a BLS aggregate signature from the
chain's validator set.
\end{definition}

\begin{table}[h]
\centering
\caption{Trust receipt types by originating chain.}
\label{tab:receipts}
\begin{tabular}{llp{6cm}}
\toprule
\textbf{Chain} & \textbf{Receipt Type} & \textbf{Meaning} \\
\midrule
K & \textsc{KeyRegistered} & New key record created \\
K & \textsc{KeyRotated} & Key rotated to successor \\
K & \textsc{CapGranted} & Capability delegated \\
K & \textsc{CapRevoked} & Capability revoked \\
T & \textsc{QuorumApproved} & Threshold vote passed \\
T & \textsc{QuorumRejected} & Threshold vote failed \\
M & \textsc{SignatureProduced} & Threshold signature emitted \\
I & \textsc{DIDRegistered} & DID document created \\
I & \textsc{CredentialIssued} & Credential hash anchored \\
I & \textsc{CredentialRevoked} & Credential revoked \\
\bottomrule
\end{tabular}
\end{table}

\subsection{Capsule Audit Trails}

A capsule maintains a local log of operations (encrypt, decrypt, share, revoke).
For each operation, the capsule computes a digest and submits it as an
\textsc{AnchorTx} to the appropriate chain. The chain includes the digest in a
receipt. The capsule stores the receipt locally alongside the operation log entry.

\begin{theorem}[Audit Completeness]\label{thm:audit}
If a capsule faithfully anchors every operation and the trust kernel satisfies
BFT safety (Section~\ref{sec:security}), then for any operation $o$ in the
capsule's log, there exists a finalized receipt $\mathcal{R}$ on-chain such that
$\mathcal{R}.\mathsf{digest} = H(o)$.
\end{theorem}

The on-chain receipt stores only the digest---never the operation payload. An
auditor can verify completeness by comparing the capsule's local log against the
receipt chain, without the chain ever learning the contents of the operations.

%% ===================================================================
\section{Provider Accountability}\label{sec:provider}
%% ===================================================================

The trust kernel enables a critical property: \emph{provider-independent
auditability}. A capsule's data may be hosted by any provider (cloud, on-premise,
edge), but the trust trail remains on-chain regardless of provider.

\subsection{Provider Switching}

When a capsule migrates from provider $P_1$ to provider $P_2$:

\begin{enumerate}
    \item The capsule owner initiates a key rotation on the K-Chain, generating
          new wrapping keys accessible only to $P_2$.
    \item $P_1$ produces a \textsc{TransferReceipt} attesting to the data
          handoff, anchored on-chain.
    \item $P_2$ confirms receipt of the encrypted data and anchors a
          \textsc{ReceiveReceipt}.
    \item The owner revokes $P_1$'s capabilities on the K-Chain.
\end{enumerate}

At no point does the chain store the capsule's data. The chain stores only the
trust decisions (key rotation, capability revocation, transfer attestation) that
make the migration auditable.

\begin{proposition}[Provider Independence]\label{prop:provider}
Given a capsule $\mathcal{C}$ with receipt chain $\{\mathcal{R}_1, \ldots,
\mathcal{R}_n\}$ and a provider migration from $P_1$ to $P_2$, the receipt
chain after migration contains sufficient information for any third-party
auditor to verify: (a)~the identity of both providers, (b)~the time of
migration, (c)~the completeness of capability revocation from $P_1$, and
(d)~the continuity of the capsule's key lineage.
\end{proposition}

%% ===================================================================
\section{Cross-Capsule Coordination}\label{sec:cross-capsule}
%% ===================================================================

Two capsules $\mathcal{C}_A$ and $\mathcal{C}_B$ interact through the trust
kernel without either capsule exposing plaintext to the other or to the chain.

\subsection{Key Exchange}

Capsule $\mathcal{C}_A$ retrieves $\mathcal{C}_B$'s public encapsulation key
$\mathsf{ek}_B$ from the K-Chain, verifies it is active and satisfies
$\mathcal{C}_A$'s trust policy, and performs ML-KEM encapsulation:
\[
    (c, K) \gets \mathsf{ML\text{-}KEM.Encaps}(\mathsf{ek}_B)
\]
The ciphertext $c$ is delivered to $\mathcal{C}_B$ off-chain. The K-Chain
records a \textsc{CapGranted} receipt noting that $\mathcal{C}_A$ was granted
encapsulation rights for $\mathcal{C}_B$'s key. The shared secret $K$ is used
to derive a symmetric session key for subsequent encrypted communication.

\subsection{Capability Delegation}

$\mathcal{C}_A$ may delegate a capability to $\mathcal{C}_B$ by submitting a
\textsc{GrantCap} transaction to the K-Chain:
\[
    \textsc{GrantCap}(\mathsf{key}_A, \mathsf{did}_B, \mathsf{cap}, \mathsf{expiry})
\]
This creates a verifiable, time-bounded delegation chain. $\mathcal{C}_B$ can
prove to any verifier that it holds capability $\mathsf{cap}$ over
$\mathsf{key}_A$ by presenting the \textsc{CapGranted} receipt.

\subsection{Threshold Approval}

For operations requiring mutual consent (e.g., a joint contract signing),
$\mathcal{C}_A$ and $\mathcal{C}_B$ create a 2-of-2 quorum on the T-Chain:
\[
    Q = (2, 2, \{\mathsf{did}_A, \mathsf{did}_B\}, \mathsf{AND}, \mathsf{timeout})
\]
Both capsules submit their encrypted votes. The T-Chain evaluates the AND
predicate. If both approve, a \textsc{GovernanceReceipt} is emitted that
authorizes the subsequent action (e.g., M-Chain signing, K-Chain rotation).

\begin{theorem}[Cross-Capsule Confidentiality]\label{thm:cross-capsule}
In the cross-capsule coordination protocol, the trust kernel learns:
(i)~the identities of the participating capsules, (ii)~the type of
coordination (key exchange, delegation, threshold approval), and
(iii)~the boolean outcome. The trust kernel does not learn the plaintext
of any data exchanged between capsules, any attribute values used in
ZK proofs, or any vote values in TFHE-encrypted quorum evaluations.
\end{theorem}

%% ===================================================================
\section{Security Model}\label{sec:security}
%% ===================================================================

\subsection{Trust Assumptions}

The trust kernel's security rests on the following assumptions:

\begin{enumerate}
    \item \textbf{BFT consensus}: each chain (K, T, M, I) is secured by a
          validator set $V$ with $|V| = n$ and the adversary controls at most
          $f < n/3$ validators. Under this assumption, the chain provides
          safety (no conflicting receipts) and liveness (receipts are eventually
          finalized).
    \item \textbf{Threshold honesty}: for $t$-of-$n$ protocols, at least $t$
          participants are honest. This is strictly weaker than the BFT
          assumption when $t \leq \lfloor n/3 \rfloor + 1$.
    \item \textbf{Cryptographic hardness}: ML-KEM is IND-CCA2 secure under the
          Module-LWE assumption; ML-DSA is EUF-CMA secure under Module-LWE and
          Module-SIS; TFHE is IND-CPA secure under LWE; FROST is EUF-CMA
          secure in the random oracle model under the discrete logarithm
          assumption; CGGMP21 is EUF-CMA secure under the strong RSA and DDH
          assumptions.
    \item \textbf{Capsule integrity}: each capsule correctly implements its
          local encryption and anchoring protocols. A compromised capsule can
          leak its own data but cannot forge receipts or compromise other
          capsules through the trust kernel.
\end{enumerate}

\subsection{BFT Threshold Analysis}

\begin{theorem}[Trust Kernel Safety]\label{thm:safety}
If fewer than $n/3$ validators on each chain are Byzantine, then for any
two receipts $\mathcal{R}_1, \mathcal{R}_2$ with the same
$(\mathsf{chain}, \mathsf{block}, \mathsf{txHash})$, it holds that
$\mathcal{R}_1 = \mathcal{R}_2$.
\end{theorem}

\begin{proof}
This follows directly from the BFT safety property of the underlying consensus
protocol. Two conflicting receipts at the same position would require two
conflicting blocks at the same height, which the BFT consensus prevents
when $f < n/3$. See the mechanized proof in \texttt{Consensus/BFT.lean}
(Section~\ref{sec:verification}).
\end{proof}

\subsection{Post-Quantum Migration Path}

The trust kernel is designed for incremental PQ migration:

\begin{enumerate}
    \item \textbf{Phase 1 (current)}: ML-KEM for key wrapping on K-Chain;
          classical signatures (Ed25519, secp256k1) for transaction signing.
          All wrapped key material is PQ-safe at rest.
    \item \textbf{Phase 2}: ML-DSA for transaction signatures on K-Chain and
          I-Chain. T-Chain and M-Chain retain classical signatures for
          threshold protocol compatibility.
    \item \textbf{Phase 3}: Ringtail (lattice-based threshold signatures)
          replaces FROST on M-Chain, providing full PQ security for threshold
          operations.
\end{enumerate}

Each phase is activated by a governance vote on the T-Chain. The K-Chain's
rotation policy (Definition~\ref{def:rotation}) includes the
$\mathsf{successor\_alg}$ field, enabling per-key migration.

%% ===================================================================
\section{Formal Verification}\label{sec:verification}
%% ===================================================================

All critical protocols in the trust kernel have corresponding Lean~4 mechanized
proofs. The proof modules are organized as follows:

\begin{table}[h]
\centering
\caption{Lean~4 proof modules for trust kernel properties.}
\label{tab:proofs}
\begin{tabular}{lp{7cm}}
\toprule
\textbf{Module} & \textbf{Property} \\
\midrule
\texttt{Consensus/BFT} & Safety and liveness under $f < n/3$ Byzantine
    validators \\
\texttt{Crypto/FROST} & Unforgeability of FROST threshold Schnorr
    signatures \\
\texttt{Crypto/Threshold/CGGMP21} & Threshold ECDSA composition security
    and identifiable abort correctness \\
\texttt{Trust/Authority} & Trust delegation chain integrity: no capability
    escalation, revocation finality \\
\bottomrule
\end{tabular}
\end{table}

\begin{remark}
The Lean~4 proofs verify the \emph{protocol logic}, not the implementation.
Implementation correctness is established through the type-safe Go and Rust
codebases (\texttt{luxfi/frost}, \texttt{luxfi/threshold}, \texttt{luxfi/lattice})
with property-based testing against the Lean~4 reference specifications.
\end{remark}

The proof of trust delegation integrity (Theorem~\ref{thm:signing-auth}) is
mechanized in \texttt{Trust/Authority.lean} as follows: the proof establishes
that the authorization graph is acyclic, that revocation propagates to all
downstream delegations, and that no delegation chain can exceed the
capabilities of its root authority.

%% ===================================================================
\section{Related Work}\label{sec:related}
%% ===================================================================

\subsection{UCAN (Fission)}

User Controlled Authorization Networks (UCANs) \cite{ucan-fission} implement
capability-based authorization using chained JWTs. Each UCAN token contains
an issuer, audience, and set of attenuated capabilities. The trust kernel
shares the capability delegation model but differs in three ways:
(1)~capabilities are registered on-chain rather than embedded in bearer tokens,
preventing capability replay after revocation;
(2)~delegation chains are verified by the chain's validator set rather than
by the relying party alone;
(3)~key rotation is enforced by policy rather than requiring token reissuance.

\subsection{Ceramic Anchors}

Ceramic Network \cite{ceramic-anchors} anchors document state on Ethereum,
providing a tamper-evident log for mutable documents. The trust kernel
generalizes this pattern: where Ceramic anchors document CIDs, the trust
kernel anchors \emph{trust decisions} (key lifecycle, governance votes,
signatures). This enables audit trails for operations that have no associated
document (e.g., a capability revocation or a threshold vote rejection).

\subsection{Keybase}

Keybase \cite{keybase} maintained a global key directory with Merkle tree
anchoring to Bitcoin. The K-Chain provides equivalent functionality---a
global key registry with on-chain anchoring---but extends it with:
(1)~policy-enforced rotation,
(2)~ML-KEM wrapping for post-quantum safety,
(3)~capability-scoped key usage rather than all-or-nothing key trust.

\subsection{Signal Protocol Key Management}

The Signal Protocol \cite{signal-protocol} uses the X3DH key agreement and
Double Ratchet for forward secrecy. The trust kernel does not replace
end-to-end encryption protocols but provides the \emph{root of trust} for
key discovery. Where Signal relies on a centralized key server, the trust
kernel provides a decentralized, auditable key registry with verifiable
rotation history.

%% ===================================================================
\section{Conclusion}\label{sec:conclusion}
%% ===================================================================

Lux is not a database. Lux is a trust kernel.

This paper has formalized the trust kernel abstraction: the minimum on-chain
substrate that makes sovereign applications composable, auditable, and
provider-independent. The four chains---K (key policy), T (threshold governance),
M (MPC signing), I (decentralized identity)---provide exactly the trust services
that local encrypted vaults cannot provide on their own: global identity,
cross-party coordination, and settlement finality.

The trust kernel stores no application data. It stores commitments to
decisions about data. Every decision produces a receipt. Receipts are
permanent, tamper-evident, and verifiable by any party. This architecture
enables capsules to switch providers, delegate capabilities to other capsules,
and participate in threshold governance---all without revealing plaintext to
the chain or to each other.

The post-quantum migration path ensures that trust decisions made today remain
verifiable under future quantum adversaries. The mechanized Lean~4 proofs
provide formal guarantees for the critical protocol properties: consensus
safety, signature unforgeability, threshold composition security, and trust
delegation integrity.

The trust kernel is the answer to a question that the blockchain industry has
been asking incorrectly. The question is not ``how do we put data on-chain.''
The question is ``what is the minimum we must put on-chain so that everything
else can stay off-chain and still be trusted.'' Four things: identity, key
policy, threshold governance, and receipts.

%% ===================================================================
%% References
%% ===================================================================

\begin{thebibliography}{99}

\bibitem{lux-capsule-architecture}
Z.~Kelling.
\newblock Lux Capsule Architecture: Sovereign Encrypted Vaults with Trust Kernel Binding.
\newblock Lux Industries, 2026.

\bibitem{lux-mchain-mpc}
Z.~Kelling and V.~Seesahai.
\newblock M-Chain: Decentralized Multi-Party Computation Custody with Quantum-Safe Threshold Signatures.
\newblock Lux Industries, 2023.

\bibitem{lux-id-did-specification}
Z.~Kelling.
\newblock Lux ID: Decentralized Identity and Universal Access Management.
\newblock Lux Industries, 2020.

\bibitem{lux-pq-crypto-suite}
Z.~Kelling.
\newblock Post-Quantum Cryptographic Suite for EVM: ML-KEM, ML-DSA, and SLH-DSA as Native Precompiles.
\newblock Lux Industries, 2024.

\bibitem{lux-tfhe}
Z.~Kelling.
\newblock Torus Threshold Fully Homomorphic Encryption: Boolean Circuits on Encrypted Data with Distributed Decryption.
\newblock Lux Industries, 2025.

\bibitem{lux-warp-messaging}
Z.~Kelling.
\newblock Lux Warp Messaging: Cross-Chain Communication Protocol.
\newblock Lux Industries, 2023.

\bibitem{ucan-fission}
B.~Frazee, P.~Willison, and B.~Holmgren.
\newblock UCAN: User Controlled Authorization Networks.
\newblock Fission Codes, 2021.

\bibitem{ceramic-anchors}
J.~Clark, M.~Zargham, and D.~Rees.
\newblock Ceramic: A Decentralized Data Network for Web3.
\newblock 3Box Labs, 2022.

\bibitem{keybase}
M.~Krohn, M.~Mazieres, and C.~Coyne.
\newblock Keybase: Public Key Crypto for Everyone, Publicly Auditable Proofs of Identity.
\newblock Keybase, Inc., 2014.

\bibitem{signal-protocol}
T.~Perrin and M.~Marlinspike.
\newblock The X3DH Key Agreement Protocol.
\newblock Signal Foundation, 2016.

\bibitem{frost}
C.~Komlo and I.~Goldberg.
\newblock FROST: Flexible Round-Optimized Schnorr Threshold Signatures.
\newblock In \emph{Selected Areas in Cryptography (SAC)}, 2020.

\bibitem{cggmp21}
R.~Canetti, R.~Gennaro, S.~Goldfeder, N.~Makriyannis, and U.~Peled.
\newblock UC Non-Interactive, Proactive, Threshold ECDSA with Identifiable Aborts.
\newblock In \emph{ACM CCS}, 2020.

\bibitem{mlkem}
NIST.
\newblock FIPS 203: Module-Lattice-Based Key-Encapsulation Mechanism Standard.
\newblock National Institute of Standards and Technology, 2024.

\bibitem{mldsa}
NIST.
\newblock FIPS 204: Module-Lattice-Based Digital Signature Standard.
\newblock National Institute of Standards and Technology, 2024.

\end{thebibliography}

\end{document}
